%%%%%%%%%%%%%%%%%%%%%%%%%%%%%%%%%%%%%%%%%%%%%%%%%%%%%%%%%%%%%%%%%%%%%%%%%%%
% Handout und Uebungsblatt zu Lehreinheit 9
% Das Interviewtranskript (Aufgabe 3) ist konstruiert.
%%%%%%%%%%%%%%%%%%%%%%%%%%%%%%%%%%%%%%%%%%%%%%%%%%%%%%%%%%%%%%%%%%%%%%%%%%%
\begin{handout}{Qualitative Methoden}

\textbf{Worum es heute ging.} Sie haben denselben Effekt zweimal erlebt,
mit identischen Handlungen. Nur der Sprechtext war anders: einmal
sachlich, einmal \enquote{magisch}. Die Zahl auf Ihrer Karte zeigt, ob
sich Ihr Erleben verschoben hat; Ihre drei Wörter zeigen, worin. Genau
dafür gibt es qualitative Forschung: Sie fragt, wie Menschen eine
Situation deuten, und behandelt die eigene Sprache als Teil der Daten.

\textbf{Kernbegriffe}
\begin{description}
  \item[Verstehen statt Messen] Qualitative Forschung fragt nach
    Bedeutungen und Begründungen, arbeitet mit wenigen, gezielt gewählten
    Fällen und bildet Kategorien oft erst am Material.
  \item[Konstruktivismus] Wirklichkeit wird in Interaktion und Sprache
    mit hervorgebracht \parencite{bergerluckmann1966}. Auch Interviewfragen
    formen die Antworten \parencite{loftuspalmer1974}.
  \item[Intersubjektivität] Ein Befund hängt nicht an der Person, die ihn
    erhoben hat: Verfahren offenlegen, zweite Codierung, Abweichungen
    begründen.
  \item[Leitfadeninterview] offene Leitfragen je Themenblock, darunter
    Nachfragen; Bau nach \enquote{sammeln -- prüfen -- sortieren --
    subsumieren} \parencite{helfferich2011qualitaet}.
  \item[Qualitative Inhaltsanalyse] regelgeleitetes Zuordnen von
    Textstellen zu Kategorien mit Definition, Ankerbeispiel und
    Kodierregel \parencite{mayring2015inhaltsanalyse}.
\end{description}

\textbf{Gütekriterien nach \textcite{lincolnguba1985naturalistic}}

\begin{center}
\small
\begin{tabular}{@{}lll@{}}
\toprule
\textbf{Kriterium} & \textbf{Frage} & \textbf{Mittel} \\
\midrule
Credibility & Ist die Rekonstruktion glaubwürdig? & Triangulation, member checking \\
Transferability & Lässt sich übertragen? & dichte Beschreibung \\
Dependability & Ist das Vorgehen nachvollziehbar? & Dokumentation \\
Confirmability & Folgt der Befund aus den Daten? & Audit Trail, Selbstreflexion \\
\bottomrule
\end{tabular}
\end{center}

\textbf{Vier Wege zu qualitativen Daten:} Leitfadeninterview,
Fokusgruppe, Beobachtung, Dokumentenanalyse.

\begin{merkbox}
Wer fragt, formt die Antwort mit. Gute qualitative Forschung macht das
sichtbar: offene Fragen, offengelegte Regeln, eine zweite Person, die
unabhängig codiert.
\end{merkbox}

\textbf{Zum Weiterlesen:} \textcite{helfferich2011qualitaet}, Kap.~4;
\textcite{mayring2015inhaltsanalyse}, Kap.~5; \textcite{doeringbortz2016},
Kap.~10.
\end{handout}

\begin{uebungsblatt}{Leitfaden, Transkript, Codes}

\aufgabe{Die Wörter von der Tafel ordnen}[in der Stunde, Plenum, 5 Min.]
Fassen Sie die Wörter der beiden Runden zu drei bis fünf Gruppen
zusammen und geben Sie jeder Gruppe einen Namen. Sie haben damit induktiv
Kategorien gebildet.
\schreiblinien{3}

\aufgabe{Interviewleitfaden}[in der Stunde und Hausaufgabe]
Entwickeln Sie einen Leitfaden mit 8--12 Fragen zu Ihrer Forschungsfrage.
Jeder Block beginnt mit einer offenen Leitfrage; Nachfragen stellen Sie
nur, wenn der Aspekt nicht von selbst kommt.

\textbf{Meine Forschungsfrage:} \rule{0.72\textwidth}{0.3pt}

\begin{center}
\small
\renewcommand{\arraystretch}{2.4}
\begin{tabular}{|p{0.16\textwidth}|p{0.38\textwidth}|p{0.36\textwidth}|}
\hline
\textbf{Block} & \textbf{Leitfrage (offen, erzählauffordernd)} & \textbf{Nachfragen} \\
\hline
Einstieg & & \\ \hline
Thema A & & \\ \hline
Thema B & & \\ \hline
Thema C & & \\ \hline
Abschluss & \enquote{Gibt es noch etwas, das Ihnen wichtig ist?} & \\ \hline
\end{tabular}
\end{center}

\textbf{Prüfen Sie jede Frage:} \quad $\square$ offen statt ja/nein \quad
$\square$ keine Doppelfrage \quad $\square$ keine Bewertung vorgegeben
\quad $\square$ kein Fachwort \quad $\square$ ich brauche die Antwort
wirklich

\textbf{Reflexion nach dem Probe-Interview:} Welche Frage war suggestiv,
welche zu offen? Wie würden Sie sie umformulieren?
\schreiblinien{4}

\aufgabe{Mini-Inhaltsanalyse}[in der Stunde, zwei Gruppen, 15 Min.]
\textbf{Konstruiertes Transkript} (für diese Übung erfunden, keine echte
Person). Kontext: Interview mit einer Studentin (B) eine Woche nach einer
Vorlesung, in der der Dozent einen Becher-Zaubertrick gezeigt und nicht
aufgelöst hat. I = interviewende Person.

\begin{small}
\begin{description}\setlength{\itemsep}{2pt}
\item[{[1]}] \textbf{I:} Erzählen Sie mal, wie Sie die Vorlesung am
  Dienstag erlebt haben. \textbf{B:} Ehrlich gesagt war ich am Anfang
  total müde, ich hatte vorher noch Labor. Dann hat er plötzlich diese
  Becher rausgeholt, und ich dachte: Was wird das jetzt?
\item[{[2]}] \textbf{B:} Und dann war die Kugel weg. Also wirklich weg.
  Ich hab wirklich gedacht, das geht nicht. Ich hab mich richtig
  aufgerichtet.
\item[{[3]}] \textbf{B:} Mein Nachbar hat sofort gesagt, die ist im
  Ärmel. Aber er hatte ja kurze Ärmel an. Wir haben die ganze Pause
  darüber geredet, wie das gehen könnte.
\item[{[4]}] \textbf{I:} Fanden Sie es auch unangenehm, dass Sie
  getäuscht wurden? \textbf{B:} Ja, unangenehm schon irgendwie. Also,
  nicht schlimm. Aber man fühlt sich halt ein bisschen vorgeführt.
\item[{[5]}] \textbf{I:} Was ist Ihnen von der Vorlesung im Kopf
  geblieben? \textbf{B:} Das mit den zwei Dritteln. Dass man wechseln
  soll. Das hab ich abends meiner Mitbewohnerin erklärt, und die hat's
  mir nicht geglaubt. Dann haben wir's mit Tassen ausprobiert.
\item[{[6]}] \textbf{B:} Ob ich das ohne den Trick auch behalten hätte?
  Keine Ahnung. Vielleicht. Ich war halt wach danach.
\item[{[7]}] \textbf{I:} Wie würden Sie den Dozenten beschreiben?
  \textbf{B:} Der ist schon gut. Man merkt, dass der das kann, das war
  sauber gemacht. Aber man weiß bei dem nie so richtig, ob er einen
  gerade reinlegt.
\item[{[8]}] \textbf{B:} Bei der Klausur frag ich mich jetzt manchmal,
  ob da auch irgendwo ein Trick drin ist. (lacht)
\item[{[9]}] \textbf{I:} Gibt es noch etwas, das Ihnen wichtig ist?
  \textbf{B:} Ich fand's gut, dass er's nicht aufgelöst hat. Sonst hätten
  wir nicht so lange drüber nachgedacht. Obwohl mich das schon nervt,
  ehrlich gesagt.
\item[{[10]}] \textbf{B:} Ich glaub, ich hab danach mehr mitgeschrieben
  als sonst. Aber das kann auch am Thema gelegen haben.
\end{description}
\end{small}

\textbf{Kategoriensystem} (deduktiv vorgegeben):
\textbf{K1} Staunen: erlebte Überraschung, Widerspruch zur Erwartung. --
\textbf{K2} Erklärungssuche: Vermutungen, Diskussion, Ausprobieren, wie
etwas funktioniert. -- \textbf{K3} Unbehagen: Misstrauen, Ärger, Gefühl,
vorgeführt zu werden. -- \textbf{K4} Lernen: Behalten, Aufmerksamkeit,
Mitschreiben, Weitergeben von Inhalten. -- \textbf{K5} Bewertung der
Lehrperson oder Darbietung. -- \textbf{K0} Sonstiges, Kontext.
\emph{Kodierregel:} jedes Segment genau einer Kategorie zuordnen; bei
Zweifel die Kategorie des Satzes mit dem stärksten Gewicht wählen und
den Zweifel notieren.

\begin{center}
\small
\renewcommand{\arraystretch}{1.5}
\begin{tabular}{|c|p{0.08\textwidth}|p{0.08\textwidth}|p{0.08\textwidth}|p{0.5\textwidth}|}
\hline
\textbf{Segm.} & \textbf{Gr.~1} & \textbf{Gr.~2} & \textbf{gleich?} & \textbf{Begründung bei Abweichung} \\ \hline
1 & & & & \\ \hline
2 & & & & \\ \hline
3 & & & & \\ \hline
4 & & & & \\ \hline
5 & & & & \\ \hline
6 & & & & \\ \hline
7 & & & & \\ \hline
8 & & & & \\ \hline
9 & & & & \\ \hline
10 & & & & \\ \hline
\end{tabular}
\end{center}
Übereinstimmung: \rule{10mm}{0.3pt} von 10 Segmenten $=$
\rule{12mm}{0.3pt}\,\%. Was sagt die Abweichung über Intersubjektivität?
Was fällt Ihnen an Segment~4 auf?
\schreiblinien{3}

\aufgabe{Tagebuch des Staunens}[Hausaufgabe]
Wann hat eine Formulierung, ein Tonfall oder ein Rahmen Ihr Erleben
verändert, obwohl sich an der Sache nichts geändert hat?

\begin{lehrhinweis}[Erwartungshorizont]
\textbf{Aufgabe~1.} Jede plausible Gruppierung ist richtig, wenn die
Gruppennamen die enthaltenen Wörter tatsächlich abdecken. Hinweis an die
Studierenden: Die Gruppennamen sind bereits eine Deutung.

\textbf{Aufgabe~2.} Mindestens drei Blöcke mit je einer offenen Leitfrage;
keine Ja/Nein-Fragen als Leitfragen; Nachfragen konkret. Typische
Umformulierung: \enquote{Hat Ihnen das geholfen?} wird zu \enquote{Was
hat sich danach für Sie verändert?}

\textbf{Aufgabe~3} (Musterlösung; Abweichungen sind erwünscht, wenn sie
begründet werden).
\emph{Eindeutig:} 2 K1 (Ankerbeispiel), 3 K2, 6 K4, 10 K4.
\emph{Mehrdeutig, absichtlich:} 1 K0 (Kontext; wer \enquote{Was wird das
jetzt?} als Neugier liest, wählt K1, die Kodierregel spricht aber dagegen,
weil noch keine Überraschung erlebt wird); 4 K3, aber die Frage gibt das
Wort \enquote{unangenehm} vor, die Antwort übernimmt es und schwächt es
gleich ab -- das Segment ist methodisch belastet und sollte markiert
werden; 5 K4 (Weitergeben des Inhalts), mit Anteil K2 (Ausprobieren mit
Tassen); 7 K5, der zweite Satz gehört zu K3; 8 K3 (Misstrauen),
wegen \enquote{(lacht)} auch als K5 lesbar; 9 K5 (\enquote{fand's gut})
gegen K3 (\enquote{nervt}). Vier Segmente sind eindeutig; wie viele der
übrigen sechs übereinstimmen, hängt an den Grenzfällen. Lernziel ist die
Einsicht, dass
die Abweichungen an den Kategoriengrenzen entstehen und dass eine
schärfere Kodierregel (z.\,B.\ für Segmente mit zwei Gefühlslagen)
die Übereinstimmung erhöht. Segment~4 zeigt außerdem, wie eine
suggestive Frage Daten erzeugt, die man nicht mehr sauber deuten kann.
\end{lehrhinweis}
\end{uebungsblatt}
